% Lösungen Einheit 2 von 4 – Rückwärts (zusammenbau v0.1)
\einheitenkopf{Einheit 2 von 4 · Rückwärts}

\erg{14}{a) vorwärts \quad b) $0{,}25 \cdot a = 2$, also $a = 8$ (Euro) \quad c) $\frac{4}{6} \cdot 1 + \frac{2}{6} \cdot b = 3$, also $4 + 2b = 18$ und $b = 7$ \quad d) $\frac{6}{6 + x} \cdot \frac{5}{5 + x} = \frac{1}{3}$ führt auf $(6 + x)(5 + x) = 90$, also $x^2 + 11x - 60 = 0$; $x_1 = 4$, $x_2 = -15$ entfällt; also $x = 4$ \quad e) $0{,}7 \cdot n > 30$, also $n > 42{,}86$; mindestens $n = 43$ \quad f) Erwartungswert 3 heißt Augensumme 18, die verdeckten Seiten tragen zusammen $= 18 - 10 = 8$; Kandidaten 1, 3, 4 und 2, 2, 4 und 2, 3, 3; genau drei verschiedene Zahlen haben nur 2, 2, 4 und 2, 3, 3; zweimal dieselbe Zahl: bei 2, 3, 3 $\frac{9 + 4 + 1}{36} = \frac{7}{18}$, bei 2, 2, 4 $\frac{16 + 1 + 1}{36} = \frac{1}{2}$; die Beschriftung 2, 3, 3 erfüllt alle drei}
\erg{15}{a) $16q^2 + 4 \cdot 2q(1 - q) + (1 - q)^2 = 4$ ergibt $9q^2 + 6q - 3 = 0$, geteilt durch 3 die Gleichung; $q = \frac{1}{3}$ ($q = -1$ entfällt)}
\erg{16}{a) je Tag im Mittel $270 : 30 = 9$ Münzen; mit $P(10) = 0{,}95 - p$: $5p + 10 \cdot (0{,}95 - p) + 40 \cdot 0{,}05 = 9$, also $p = P(5) = 0{,}5$ und $P(10) = 0{,}45$}
\erg{17}{a) $P(\text{Sachpreis}) = 1 - 0{,}002 - 0{,}01 = 0{,}988$; $50 \cdot 0{,}002 + 10 \cdot 0{,}01 + x = 0{,}95$ (alles in Euro), also $x = 0{,}75$ (Euro), 75 Cent}
\erg{18}{a) $0{,}3a = 0{,}2b$, also $\frac{a}{b} = \frac{2}{3}$, $a : b = 2 : 3$}
\erg{19}{a) $E(k) = 2 \cdot \frac{k}{10} \cdot \frac{10 - k}{10}$ ist eine nach unten geöffnete Parabel mit Scheitel bei $k = 5$; $E(5) = 0{,}5$ (Euro)}
\erg{20}{a) Die erwartete Summe zweier Kugeln ist das Doppelte des mittleren Kugelwerts, der mittlere Wert ist also 3; da 1 und 2 darunter liegen, muss $c$ darüber liegen.}
\erg{21}{a) $n(9 - n) = (n + 1)(8 - n)$ ergibt $9n = 7n + 8$, also $n = 4$; die Erwartungswerte für 4 und 5 rote Kugeln sind gleich, die Aussage stimmt}
\erg{22}{a) Fehler: Tim setzt die Gewinnerwartung gleich dem Einsatz; fair heißt Gewinnerwartung null: $0{,}25a - 2 = 0$, also $a = 8$ (Euro) \quad b) Ausgleich auf lange Sicht heißt, dass die mittlere Auszahlung je Spiel genau der Einsatz ist; der Erwartungswert der Auszahlung wird dem Einsatz gleichgesetzt, und das ist eine Gleichung für die Unbekannte. \quad c) mittlere Erstattung $= 250 \cdot 0{,}08 = 20$ (Euro); Prämie mindestens $20 + 12 = 32$ (Euro)}
